\section{Introduction}
\label{sec:intro}
\textbf{This project investigates how a Graph Neural Network (GNN)-based music knowledge graph, combining musical features (e.g., pitch statistics and tempo) as node embeddings, song metadata (such as artists and genres) as graph structure, and user interaction history as supervision, can be used to predict user-specific music preferences.}

Although many music recommendation algorithms already exist and are widely used, most of them rely primarily on user interaction data and similarity measures between users. While suggestions based on mutual listening history overlap can be quite effective, this approach remains inherently limited and strongly favors music and artists that are already popular. Musical preference is highly subjective, and relying solely on behavioral patterns may overlook important aspects of the music itself.

In reality, musical pieces are a rich source of information, differing from one another across multiple dimensions. Each piece can be characterized by a wide range of features, from more direct attributes such as the artist or genre to strictly musical properties derived through analysis, such as pitch statistics, tempo, or chord progressions. Incorporating these musical features into the recommendation process may provide a richer and more structured representation of songs, enabling models to capture deeper relationships between musical content and user preferences.

Because music is both an art form and a structured domain, it is particularly well-suited for computational analysis and representation as a knowledge graph. By leveraging a knowledge graph in combination with Graph Neural Networks (GNNs), this project aims to model complex, multi-relational dependencies and explore their potential for innovative,  more accurate and context-aware music preference prediction.

